\section*{Chapter \ref{ch:g8:jobs-wages-unemployment} --- Jobs, Wages and Unemployment}

\begin{solution}{exo:g8:jobs-wages-unemployment:1}
About $26\,\%$ ($26.1\,\%$), in 2013.
\end{solution}

\begin{solution}{exo:g8:jobs-wages-unemployment:2}
\omterm{def:g8:jobs-wages-unemployment:labour-force}{Labour force} $18\,800 + 1\,200 = 20\,000$. \omterm{def:g8:jobs-wages-unemployment:unemployment-rate}{Unemployment rate} $1\,200 \div 20\,000 \times 100 = 6.0\,\%$.
\end{solution}

\begin{solution}{exo:g8:jobs-wages-unemployment:3}
$24\,000 \div 40\,000 \times 100 = 60\,\%$.
\end{solution}

\begin{solution}{exo:g8:jobs-wages-unemployment:4}
About 1996 (Germany $8.8\,\%$, the United Kingdom $8.2\,\%$), and back below it about 2010 (Germany
$7.0\,\%$, the United Kingdom $7.8\,\%$).
\end{solution}

\begin{solution}{exo:g8:jobs-wages-unemployment:5}
$4\,376 - 4\,122 = 254$ thousand unemployed; $254 \div 4\,376 \times 100 \approx 5.8\,\%$.
\end{solution}

\begin{solution}{exo:g8:jobs-wages-unemployment:6}
Unemployed $1\,100$, \omterm{def:g8:jobs-wages-unemployment:labour-force}{labour force} $19\,900$: $1\,100 \div 19\,900 \times 100 \approx 5.5\,\%$. The
\omterm{def:g8:jobs-wages-unemployment:participation-rate}{participation rate} falls: the \omterm{def:g8:jobs-wages-unemployment:labour-force}{labour force} is smaller and the population of working age has not
changed.
\end{solution}

\begin{solution}{exo:g8:jobs-wages-unemployment:7}
\omterm{def:g8:inflation-and-price-indices:real-wage}{Real wages} by $14.2$ (index $1\,421.6$ with 1913 = 100); real \omterm{def:g8:gdp:gdp-per-person}{GDP per person} by $9.4$.
\end{solution}

\begin{solution}{exo:g8:jobs-wages-unemployment:8}
At $2\,\%$ a year, \omterm{def:g7:specialisation:output-per-worker}{output per worker} doubles in about $70 \div 2 = 35$ years. If the model holds, \omterm{def:g8:inflation-and-price-indices:real-wage}{real
wages} roughly double too.
\end{solution}

\begin{solution}{exo:g8:jobs-wages-unemployment:9}
For example: in the United States, from 1913 to 2020, \omterm{def:g8:inflation-and-price-indices:real-wage}{real wages} were multiplied by $4.6$ and real \omterm{def:g8:gdp:gdp-per-person}{GDP
per person} by $6.3$; in Sweden, \omterm{def:g8:inflation-and-price-indices:real-wage}{real wages} by $14.2$ and real \omterm{def:g8:gdp:gdp-per-person}{GDP per person} by $9.4$. The two did not
rise by the same factor.
\end{solution}

\begin{solution}{exo:g8:jobs-wages-unemployment:10}
$(300 \times 3\,000 + 100 \times 7\,000) \div 400 = (900\,000 + 700\,000) \div 400 = 4\,000$ dalens a
month.
\end{solution}

\begin{solution}{exo:g8:jobs-wages-unemployment:11}
(a) \omterm{def:g8:jobs-wages-unemployment:labour-force}{Labour force} $30\,000$; \omterm{def:g8:jobs-wages-unemployment:unemployment-rate}{unemployment rate} $1\,500 \div 30\,000 \times 100 = 5.0\,\%$;
\omterm{def:g8:jobs-wages-unemployment:participation-rate}{participation rate} $30\,000 \div 50\,000 \times 100 = 60.0\,\%$. (b) Unemployed $1\,000$, \omterm{def:g8:jobs-wages-unemployment:labour-force}{labour force}
$29\,500$: \omterm{def:g8:jobs-wages-unemployment:unemployment-rate}{unemployment rate} $1\,000 \div 29\,500 \times 100 \approx 3.4\,\%$; \omterm{def:g8:jobs-wages-unemployment:participation-rate}{participation rate}
$29\,500 \div 50\,000 \times 100 = 59.0\,\%$. (c) No: the \omterm{def:g8:jobs-wages-unemployment:unemployment-rate}{unemployment rate} fell only because 500
people left the \omterm{def:g8:jobs-wages-unemployment:labour-force}{labour force}.
\end{solution}

\begin{solution}{exo:g8:jobs-wages-unemployment:12}
For example: how many hours people work (a person working 5 hours a week who wants 35 counts as
employed); how much jobs pay and how safe or stable they are; how many people would like a job but
have stopped looking.
\end{solution}

\begin{solution}{pb:g8:jobs-wages-unemployment:1}
\textbf{1.} Under 15: 6, 15. Employed: 1, 5, 8, 10, 11, 14, 16, 18, 20. Unemployed: 4, 9, 12, 19.
\omterm{def:g8:jobs-wages-unemployment:outside}{Outside the labour force}: 2, 3, 7, 13, 17.
\textbf{2.} 2 under 15, 9 employed, 4 unemployed, 5 \omterm{def:g8:jobs-wages-unemployment:outside}{outside the labour force} ($2 + 9 + 4 + 5 = 20$).
\textbf{3.} She has stopped looking for work.
\textbf{4.} He is not available: he cannot start work before June, months after the survey week.
\textbf{5.} $9 + 4 = 13$.
\textbf{6.} $4 \div 13 \times 100 \approx 30.8\,\%$.
\textbf{7.} Working age: $20 - 2 = 18$; $13 \div 18 \times 100 \approx 72.2\,\%$.
\textbf{8.} Unemployed 5, \omterm{def:g8:jobs-wages-unemployment:labour-force}{labour force} 14: \omterm{def:g8:jobs-wages-unemployment:unemployment-rate}{unemployment rate} $5 \div 14 \times 100 \approx 35.7\,\%$;
\omterm{def:g8:jobs-wages-unemployment:participation-rate}{participation rate} $14 \div 18 \times 100 \approx 77.8\,\%$.
\textbf{9.} 2000: $240 \div 4\,000 \times 100 = 6.0\,\%$. 2025: $279 \div 4\,645 \times 100 \approx
6.0\,\%$.
\textbf{10.} 2008: $290 \div 4\,196 \times 100 \approx 6.9\,\%$. 2009: $338 \div 4\,221 \times 100
\approx 8.0\,\%$. $338 - 290 = 48$ thousand more unemployed.
\textbf{11.} $4\,645 \div 7\,801 \times 100 \approx 59.5\,\%$.
\textbf{12.} 2012: $279 \div 4\,298 \times 100 \approx 6.5\,\%$; 2020: $397 \div 4\,508 \times 100
\approx 8.8\,\%$; 2021: $218 \div 4\,535 \times 100 \approx 4.8\,\%$. The highest is 2020, with
$8.8\,\%$.
\end{solution}
